% =============================================================================
%  Terminish École — Manuel du chef d'établissement
%  Créer une classe, inscrire les élèves, créer les comptes des enseignants,
%  programmer un examen blanc, saisir les notes, lire les résultats.
%
%  Compilation :
%    - Overleaf : compilateur pdfLaTeX (par défaut suffit)
%    - en local  : pdflatex manuel.tex  (deux fois, pour le sommaire)
%
%  Auteur : Terminish École — version 1.1 — octobre 2026
%  Établissement de référence : CEG BOROME — HIHEATRO (Togo)
%  Adresse de l'application : https://kodjone.alwaysdata.net
% =============================================================================
\documentclass[11pt,a4paper]{article}

% ---------------------------------------------------------------- encodage
\usepackage[utf8]{inputenc}
\usepackage[T1]{fontenc}
\usepackage{lmodern}
\usepackage[french]{babel}
\frenchspacing

% ---------------------------------------------------------------- mise en page
\usepackage[a4paper,top=2.2cm,bottom=2cm,left=2cm,right=2cm]{geometry}
\usepackage{graphicx}
\usepackage{xcolor}
\usepackage{amssymb}
\usepackage{array}
\usepackage{tabularx}
\usepackage{booktabs}
\usepackage{enumitem}
\usepackage{fancyhdr}
\usepackage{titlesec}
\usepackage{needspace}
\usepackage[most]{tcolorbox}
\usepackage[unicode,hidelinks]{hyperref}

% ---------------------------------------------------------------- couleurs
\definecolor{vert}{HTML}{0B6B4A}
\definecolor{vertmoyen}{HTML}{1F9D6B}
\definecolor{vertclair}{HTML}{E8F5EF}
\definecolor{or}{HTML}{B07A10}
\definecolor{orclair}{HTML}{FDF6E3}
\definecolor{rouge}{HTML}{B3261E}
\definecolor{rougeclair}{HTML}{FDECEB}
\definecolor{bleu}{HTML}{1D4E89}
\definecolor{bleuclair}{HTML}{EEF4FB}
\definecolor{gris}{HTML}{5B6B66}
\definecolor{grisclair}{HTML}{F4F6F5}
\definecolor{bord}{HTML}{D7E0DC}
\definecolor{encre}{HTML}{16241F}

% ---------------------------------------------------------------- titres
\titleformat{\section}
  {\needspace{6\baselineskip}\normalfont\Large\bfseries\color{vert}}
  {\thesection.}{0.6em}{}
\titlespacing*{\section}{0pt}{1.6\baselineskip}{0.5\baselineskip}

\titleformat{\subsection}
  {\needspace{4\baselineskip}\normalfont\large\bfseries\color{encre}}
  {\thesubsection}{0.6em}{}
\titlespacing*{\subsection}{0pt}{1.1\baselineskip}{0.35\baselineskip}

\titleformat{\subsubsection}
  {\normalfont\normalsize\bfseries\color{gris}}
  {}{0em}{}

% ---------------------------------------------------------------- en-tête / pied
\pagestyle{fancy}
\fancyhf{}
\renewcommand{\headrulewidth}{0.4pt}
\renewcommand{\footrulewidth}{0.4pt}
\fancyhead[L]{\footnotesize\color{gris}Terminish École — Manuel du chef d'établissement}
\fancyhead[R]{\footnotesize\color{gris}CEG BOROME — HIHEATRO}
\fancyfoot[L]{\footnotesize\color{gris}Version 1.1 — octobre 2026}
\fancyfoot[R]{\footnotesize\bfseries\color{vert}\thepage}

% ---------------------------------------------------------------- listes
\setlist[itemize]{leftmargin=1.3em,itemsep=2pt,topsep=3pt}
\setlist[enumerate]{leftmargin=1.6em,itemsep=2pt,topsep=3pt}

% ---------------------------------------------------------------- boîtes
\newtcolorbox{boitecle}[1][]{
  enhanced, breakable, colback=vertclair, colframe=vert,
  boxrule=0pt, leftrule=2.6pt, arc=1.2pt, outer arc=1.2pt,
  left=8pt, right=8pt, top=6pt, bottom=6pt,
  fonttitle=\bfseries\small, coltitle=vert,
  title={\textcolor{vert}{\checkmark}\enspace #1}
}
\newtcolorbox{boiteastuce}[1][]{
  enhanced, breakable, colback=vertclair, colframe=vertmoyen,
  boxrule=0pt, leftrule=2.6pt, arc=1.2pt, outer arc=1.2pt,
  left=8pt, right=8pt, top=6pt, bottom=6pt,
  fonttitle=\bfseries\small, coltitle=vert,
  title={\textcolor{vertmoyen}{$\star$}\enspace #1}
}
\newtcolorbox{boiteattention}[1][]{
  enhanced, breakable, colback=orclair, colframe=or,
  boxrule=0pt, leftrule=2.6pt, arc=1.2pt, outer arc=1.2pt,
  left=8pt, right=8pt, top=6pt, bottom=6pt,
  fonttitle=\bfseries\small, coltitle=or,
  title={\textcolor{or}{!}\enspace #1}
}
\newtcolorbox{boiteerreur}[1][]{
  enhanced, breakable, colback=rougeclair, colframe=rouge,
  boxrule=0pt, leftrule=2.6pt, arc=1.2pt, outer arc=1.2pt,
  left=8pt, right=8pt, top=6pt, bottom=6pt,
  fonttitle=\bfseries\small, coltitle=rouge,
  title={\textcolor{rouge}{$\times$}\enspace #1}
}
\newtcolorbox{boiteinfo}[1][]{
  enhanced, breakable, colback=bleuclair, colframe=bleu,
  boxrule=0pt, leftrule=2.6pt, arc=1.2pt, outer arc=1.2pt,
  left=8pt, right=8pt, top=6pt, bottom=6pt,
  fonttitle=\bfseries\small, coltitle=bleu,
  title={\textcolor{bleu}{i}\enspace #1}
}

% boîte pour les commandes et les formats à recopier
\newtcolorbox{boitecode}{
  enhanced, breakable, colback=grisclair, colframe=bord,
  boxrule=0.4pt, arc=1.2pt, left=8pt, right=8pt, top=6pt, bottom=6pt,
  fontupper=\ttfamily\small
}

% ---------------------------------------------------------------- raccourcis
\newcommand{\menu}[1]{\textbf{«\,#1\,»}}
\newcommand{\bouton}[1]{\textbf{#1}}
\newcommand{\champ}[1]{\textit{#1}}
\newcommand{\coeur}{\textcolor{vert}{\checkmark}}
\newcommand{\croix}{\textcolor{rouge}{$\times$}}

\title{Manuel du chef d'établissement}
\author{Terminish École}
\date{Version 1.1 — octobre 2026}

% =============================================================================
\begin{document}

% ---------------------------------------------------------------- couverture
\begin{titlepage}
\pagecolor{vert}
\color{white}
\vspace*{3.2cm}
{\small\color{vertclair}TERMINISH ÉCOLE}\\[2mm]
{\Huge\bfseries Manuel du chef\\[1mm] d'établissement}\\[6mm]
{\large\color{vertclair}Classes · Élèves · Enseignants · Examens blancs · Notes · Résultats}\\[3cm]
{\large CEG BOROME — HIHEATRO (Togo)}\\[2mm]
{\large https://kodjone.alwaysdata.net}\\[2mm]
{\color{vertclair}Version 1.1 — octobre 2026}
\vfill
{\small\color{vertclair}Ce manuel décrit exactement ce que l'application affiche : chaque champ,\\
chaque bouton, chaque message d'erreur y figure tel quel.}
\end{titlepage}
\pagecolor{white}
\color{encre}

% ---------------------------------------------------------------- sommaire
\tableofcontents
\newpage

% =============================================================================
\section{Comment utiliser ce manuel}

Ce manuel ne porte que sur le travail du \textbf{chef d'établissement} : c'est lui
qui crée les classes, inscrit les élèves, ouvre les comptes des enseignants,
programme les examens blancs, puis lit et publie les résultats. Les enseignants,
eux, n'ont qu'un écran à connaître : la saisie des notes.

\begin{boiteastuce}[Le principe à retenir]
L'ordre des opérations n'est pas une recommandation : c'est une
\textbf{contrainte du logiciel}. Le site refuse de créer une classe s'il n'y a
aucune matière, et refuse d'affecter un enseignant à une matière qui n'est pas
enseignée dans la classe. Respectez la chaîne ci-dessous, et vous ne verrez
aucun message de blocage.
\end{boiteastuce}

\subsection{L'ordre obligatoire des six étapes}

\begin{center}
\renewcommand{\arraystretch}{1.25}
\begin{tabularx}{\textwidth}{@{}c >{\raggedright}p{3.9cm} >{\raggedright}p{2.9cm} X@{}}
\toprule
\textbf{\#} & \textbf{Étape} & \textbf{Menu} & \textbf{Pourquoi avant la suivante} \\
\midrule
1 & Créer les \textbf{matières} & Matières & La classe choisit ses matières dans cette liste \\
2 & Créer les \textbf{classes} et choisir les coefficients & Classes & L'élève doit avoir une classe d'accueil \\
3 & Créer les comptes des \textbf{enseignants} & Enseignants & On affecte un professeur à une classe \emph{et} une matière \\
4 & \textbf{Inscrire les élèves} & Élèves & — \\
5 & \textbf{Saisir les notes} & Saisie des notes & Il faut classe + matière + période \\
6 & Lire et publier les \textbf{résultats} & Bulletins & Il faut des notes pour calculer \\
\bottomrule
\end{tabularx}
\end{center}

\begin{boiteattention}[Le message que vous verrez si l'ordre n'est pas respecté]
Si vous cliquez sur \bouton{+ Nouvelle classe} alors qu'aucune matière n'existe,
le site affiche : \menu{Créez d'abord au moins une matière}. Ce n'est pas une
panne : c'est la première étape qui manque.
\end{boiteattention}

\subsection{Les quatre rôles}

\begin{center}
\renewcommand{\arraystretch}{1.25}
\begin{tabularx}{\textwidth}{@{} >{\raggedright}p{3.4cm} >{\raggedright}p{4.2cm} X@{}}
\toprule
\textbf{Rôle} & \textbf{Ce qu'il voit} & \textbf{Ce qu'il peut faire} \\
\midrule
Chef d'établissement & Tous les écrans, toutes les classes & Tout : réglages, classes, élèves,
comptes des enseignants, publication, abonnement \\
Enseignant & Ses classes et ses matières seulement & Saisir et corriger ses notes ;
ne gère aucun identifiant \\
Élève & Ses résultats publiés & Consulter ses bulletins, une fois publiés \\
Éditeur de la plateforme & La console commerciale & Créer les établissements clients ;
n'accède pas aux notes \\
\bottomrule
\end{tabularx}
\end{center}

\begin{boitecle}[Une règle simple]
Seul le chef d'établissement crée, modifie et supprime les comptes des
enseignants. Un professeur ne peut pas s'inscrire lui-même : c'est volontaire,
pour qu'aucun compte fantôme ne puisse saisir des notes.
\end{boitecle}

\newpage

% =============================================================================
\section{Créer une classe}

\textbf{Chemin :} menu de gauche \menu{Classes} $\rightarrow$ bouton bleu
\bouton{+ Nouvelle classe}.

La page présente deux groupes : \textbf{Collège} (trimestres T1, T2, T3) et
\textbf{Lycée} (semestres S1, S2). Une classe appartient à l'un des deux, et cela
décide de tout le reste.

\subsection{Les sept champs du formulaire}

\begin{center}
\renewcommand{\arraystretch}{1.25}
\begin{tabularx}{\textwidth}{@{} >{\raggedright}p{3.5cm} c >{\raggedright}p{3.4cm} X@{}}
\toprule
\textbf{Champ} & \textbf{Obligatoire} & \textbf{Exemple} & \textbf{Effet} \\
\midrule
Nom de la classe & oui & 6ème A & Nom imprimé sur le bulletin \\
Niveau & non & 6ème & Sert aux statistiques par niveau \\
Cycle & oui & Collège ou Lycée & \textbf{Décide des périodes} : T1/T2/T3 ou S1/S2 \\
Série & non & A, C, D & Lycée uniquement ; vide au collège \\
Professeur titulaire & non & M. ADJOVI & Il \textbf{signe} les bulletins de cette classe \\
Nom affiché si aucun compte & non & M. ADJOVI & Pour un professeur sans compte sur le site \\
Matières et coefficients & oui & voir ci-dessous & Seules les matières cochées paraissent au bulletin \\
\bottomrule
\end{tabularx}
\end{center}

\begin{boiteattention}[Le cycle est le choix le plus lourd de conséquences]
Une classe créée en \textbf{Collège} travaillera sur T1, T2, T3 ; une classe créée
en \textbf{Lycée} sur S1, S2. Le cycle reste modifiable ensuite par le crayon,
mais les notes déjà saisies sur les périodes de l'ancien cycle ne suivent pas.
Vérifiez ce champ avant de valider.
\end{boiteattention}

\subsection{Le tableau « Matières et coefficients »}

C'est la partie décisive du formulaire. Il affiche une ligne par matière créée,
avec trois colonnes : le nom, la case \menu{Enseignée}, le coefficient.

\begin{center}
\renewcommand{\arraystretch}{1.3}
\begin{tabular}{@{}lcc@{}}
\toprule
\textbf{Matière} & \textbf{Enseignée} & \textbf{Coefficient} \\
\midrule
Français & $\blacksquare$ & 3 \\
Mathématiques & $\blacksquare$ & 4 \\
Anglais & $\blacksquare$ & 2 \\
Physique-Chimie et Technologie (PCT) & $\blacksquare$ & 2 \\
Sciences de la Vie et de la Terre & $\blacksquare$ & 1 \\
Histoire-Géographie & $\blacksquare$ & 2 \\
Éducation civique & $\blacksquare$ & 1 \\
Éducation physique et sportive & $\blacksquare$ & 1 \\
Espagnol & $\square$ & 1 \\
\bottomrule
\end{tabular}
\end{center}

\begin{itemize}
  \item Case \textbf{cochée} : la matière est évaluée, elle sera \textbf{imprimée}
        sur le bulletin et sa note comptera dans la moyenne générale.
  \item Case \textbf{décochée} : la matière n'existe pas dans cette classe ;
        elle n'apparaît nulle part, ni à la saisie, ni au bulletin.
  \item \textbf{Coefficient} : son poids dans la moyenne générale, de 0,5 en 0,5
        (minimum 0,5 ; maximum 20).
\end{itemize}

\begin{boiteastuce}[Le coefficient en une phrase]
La moyenne générale n'est pas la moyenne simple des matières : chaque matière
pèse selon son coefficient. Un 12 en mathématiques (coefficient 4) pèse deux fois
plus qu'un 12 en anglais (coefficient 2).
\end{boiteastuce}

\subsection{Ce que produit la création}

Une \textbf{carte de classe} apparaît, avec le nom, l'effectif, le nom du
professeur titulaire, des pastilles de matières (\texttt{MATH $\times$4},
\texttt{FRAN $\times$3}, etc.) et quatre boutons :

\begin{center}
\renewcommand{\arraystretch}{1.25}
\begin{tabularx}{\textwidth}{@{} >{\raggedright}p{2.3cm} X@{}}
\toprule
\textbf{Bouton} & \textbf{Effet} \\
\midrule
Élèves & Ouvre la liste des élèves filtrée sur cette classe \\
Saisir & Ouvre la grille de notes filtrée sur cette classe \\
Crayon & Modifier la classe (champs, matières, coefficients, titulaire) \\
Corbeille & Supprimer la classe (refusé si des élèves y sont encore inscrits) \\
\bottomrule
\end{tabularx}
\end{center}

\subsection{Les trois erreurs classiques}

\begin{boiteerreur}[À éviter absolument]
\begin{enumerate}
  \item \textbf{Se tromper de cycle.} Une classe de collège créée en « Lycée »
        proposera S1 et S2 au lieu de T1, T2, T3.
  \item \textbf{Oublier de cocher une matière.} Elle sera absente du bulletin et
        sa note ne comptera pas : l'élève sera classé sur un total incomplet.
  \item \textbf{Écrire deux fois le même nom différemment.} \texttt{6ème A} et
        \texttt{6eme A} sont deux classes distinctes. Choisissez une orthographe
        et gardez-la partout, y compris dans l'import des élèves.
\end{enumerate}
\end{boiteerreur}

\newpage

% =============================================================================
\section{Enregistrer les élèves}

Deux méthodes existent : l'inscription individuelle, et l'import d'une liste
entière. Pour 143 élèves, l'import prend deux minutes ; l'inscription
individuelle prend une soirée.

\subsection{Un élève à la fois}

\textbf{Chemin :} \menu{Élèves} $\rightarrow$ \bouton{+ Nouvel élève}.

\begin{center}
\renewcommand{\arraystretch}{1.25}
\begin{tabularx}{\textwidth}{@{} >{\raggedright}p{4.3cm} c X@{}}
\toprule
\textbf{Champ} & \textbf{Obligatoire} & \textbf{Remarque} \\
\midrule
Nom & oui & Mis en majuscules automatiquement à l'enregistrement \\
Prénom(s) & oui & \\
Sexe & non & Masculin ou Féminin ; sert aux statistiques \\
Classe & non & Si vide, l'élève reste « Non affecté » et n'apparaît pas à la saisie \\
Date de naissance & non & Sélecteur de date, format année-mois-jour \\
Lieu de naissance & non & \\
Nom du tuteur & non & Administratif : \textbf{jamais imprimé} sur le bulletin \\
Téléphone du tuteur & non & Administratif \\
Matricule & non & \textbf{Laisser vide} : généré sous la forme année + numéro \\
\bottomrule
\end{tabularx}
\end{center}

À l'enregistrement, un message confirme et donne le matricule attribué, par
exemple : \menu{Élève inscrit — matricule 20260001}.

\begin{boiteattention}[Le tuteur ne figure jamais sur le bulletin]
Le nom et le téléphone du tuteur sont des informations de secrétariat,
conservées dans le dossier de l'élève. Le bulletin est remis à la famille : il ne
contient ni ces coordonnées, ni la mention d'un tuteur.
\end{boiteattention}

\subsection{Importer toute une liste}

\textbf{Chemin :} \menu{Élèves} $\rightarrow$ \bouton{Importer}.

Le format attendu est une ligne par élève, dans cet ordre exact :

\begin{boitecode}
NOM ; PRÉNOM ; SEXE ; DATE DE NAISSANCE ; CLASSE ; TUTEUR ; TÉLÉPHONE\\[4pt]
\hspace*{1em}ADJOVI;Ama;F;2012-03-14;6ème A;ADJOVI Kossi;+228 90 00 00 00\\
\hspace*{1em}KOUASSI;Elom;M;2011-11-02;5ème A;;
\end{boitecode}

\begin{itemize}
  \item Le séparateur peut être le point-virgule, la virgule ou une
        \textbf{tabulation} : un copier-coller direct depuis Excel fonctionne.
  \item La colonne \texttt{CLASSE} doit correspondre \textbf{exactement} à un nom
        de classe existant.
  \item Si la colonne classe est vide, utilisez le menu
        \champ{Classe par défaut} du formulaire d'import.
  \item Les matricules se génèrent tout seuls : ne les écrivez pas.
  \item Un bilan s'affiche à la fin : nombre d'élèves créés, et lignes refusées
        avec leur numéro.
\end{itemize}

\begin{boiteastuce}[Créer les classes d'abord, importer ensuite]
L'import se fait toujours après la création des classes. Si une classe manque,
laissez la colonne vide et choisissez une classe par défaut : vous déplacerez les
élèves ensuite, individuellement.
\end{boiteastuce}

\subsection{Retrouver, exporter, archiver}

\begin{itemize}
  \item La \textbf{recherche} filtre par nom, prénom ou matricule.
  \item Le filtre \textbf{Classe} isole une division ; le filtre \textbf{Statut}
        sépare les actifs des archivés.
  \item \bouton{Export CSV} télécharge la liste complète, ouvrable dans Excel :
        c'est votre sauvegarde des inscriptions.
  \item La \textbf{corbeille} d'une ligne n'efface rien : elle \textbf{archive}.
        L'élève sort des listes actives, de la saisie et des bulletins, mais son
        dossier et ses notes restent. Pour le réintégrer, repassez son statut sur
        « Actif ».
\end{itemize}

\newpage

% =============================================================================
\section{Créer le compte d'un enseignant}

\textbf{Chemin :} \menu{Enseignants} $\rightarrow$ \bouton{+ Ajouter un enseignant}.

\subsection{Les cinq zones du formulaire}

\begin{center}
\renewcommand{\arraystretch}{1.25}
\begin{tabularx}{\textwidth}{@{} >{\raggedright}p{4.2cm} c X@{}}
\toprule
\textbf{Zone} & \textbf{Obligatoire} & \textbf{Détail} \\
\midrule
Nom complet & oui & Par exemple : M. Komlan AGBEKO \\
Téléphone & non & Pour le joindre ; ne sert pas à la connexion \\
Email de connexion & oui & C'est son \textbf{identifiant}. Non modifiable ensuite : le champ
est grisé à la réédition \\
Mot de passe & oui & 6 caractères minimum ; affiché en clair pendant la saisie \\
Signature de l'enseignant & non & Photo de sa signature ; elle sera apposée sur les bulletins \\
\bottomrule
\end{tabularx}
\end{center}

\begin{boiteattention}[Notez le mot de passe avant de fermer la fenêtre]
Le mot de passe s'affiche en clair dans le formulaire, mais il n'est plus
consultable après l'enregistrement. Notez-le, puis dictez-le de vive voix à
l'intéressé. S'il est perdu : bouton \bouton{clé} de sa ligne, qui en génère un
nouveau, affiché une seule fois.
\end{boiteattention}

\subsection{La grille « Classes et matières affectées »}

C'est un tableau croisé : les \textbf{classes} en lignes, les \textbf{matières} en
colonnes. Chaque case cochée crée une affectation.

\begin{center}
\renewcommand{\arraystretch}{1.3}
\begin{tabular}{@{}lcccc@{}}
\toprule
\textbf{Classe} & \textbf{Français} & \textbf{Mathématiques} & \textbf{Anglais} & \textbf{PCT} \\
\midrule
6ème A & $\square$ & $\blacksquare$ & $\square$ & $\blacksquare$ \\
5ème B & $\square$ & $\square$ & $\square$ & $\square$ \\
\bottomrule
\end{tabular}
\end{center}

\begin{itemize}
  \item Une case \textbf{grise} signifie que la matière n'est pas enseignée dans
        cette classe : elle est impossible à cocher tant que la matière n'a pas
        été ajoutée à la classe.
  \item Chaque case cochée donne au professeur l'accès à cette classe et à cette
        matière dans l'écran \menu{Saisie des notes}.
\end{itemize}

\subsection{Les quatre boutons de la ligne}

\begin{center}
\renewcommand{\arraystretch}{1.25}
\begin{tabularx}{\textwidth}{@{} >{\raggedright}p{2.6cm} X@{}}
\toprule
\textbf{Bouton} & \textbf{Effet} \\
\midrule
Crayon & Modifier nom, téléphone, mot de passe, signature, affectations. L'email reste figé \\
Clé & \textbf{Nouveaux identifiants} : génère un mot de passe, affiché une seule fois, et
ferme toutes ses sessions en cours \\
Interdit & Retire l'accès immédiatement (sessions fermées, affectations retirées) \\
Validation & Rétablit l'accès \\
\bottomrule
\end{tabularx}
\end{center}

Les demandes de mot de passe envoyées par les professeurs apparaissent en haut de
cette page, dans un bloc orangé \menu{Demandes de mot de passe}, avec le bouton
\bouton{Nouveaux identifiants} et le bouton de traitement.

\begin{boiteerreur}[L'erreur numéro un avec les enseignants]
\textbf{Créer le compte ne suffit pas.} Un enseignant dont \emph{aucune} case n'est
cochée dans la grille d'affectations se connecte, puis ne voit rien : l'écran
affiche \menu{Aucune classe ne vous est affectée}. Vérifiez toujours la colonne
\menu{Affectations} de la liste des enseignants : elle doit montrer au moins une
pastille du type \texttt{6ème A · Mathématiques}.
\end{boiteerreur}

\newpage

% =============================================================================
\section{Programmer un examen blanc}

Un examen blanc se compose en deux temps : \textbf{un examen} (le contenant :
nom, dates), puis \textbf{des épreuves} (le contenu : une matière, une classe, une
date, une durée, un coefficient). Un examen sans épreuve ne produit aucun
résultat.

\subsection{Créer l'examen}

\textbf{Chemin :} \menu{Examens blancs} $\rightarrow$ \bouton{+ Nouvel examen blanc}.

\begin{center}
\renewcommand{\arraystretch}{1.25}
\begin{tabularx}{\textwidth}{@{} >{\raggedright}p{4.4cm} X@{}}
\toprule
\textbf{Champ} & \textbf{Exemple} \\
\midrule
Nom de l'examen & Examen blanc n°1 — Tle D \\
Période de rattachement & 2ème Trimestre, ou « Aucune » \\
Début & jour de la première épreuve \\
Fin & jour de la dernière épreuve \\
\bottomrule
\end{tabularx}
\end{center}

\subsection{Ajouter les épreuves}

Sur la carte de l'examen : \bouton{+ Ajouter une épreuve}.

\begin{center}
\renewcommand{\arraystretch}{1.25}
\begin{tabularx}{\textwidth}{@{} >{\raggedright}p{3.6cm} c X@{}}
\toprule
\textbf{Champ} & \textbf{Défaut} & \textbf{Rôle} \\
\midrule
Classe & — & L'épreuve concerne une seule classe \\
Matière & — & La liste ne propose que les matières de la classe choisie \\
Date & aujourd'hui & Jour de l'épreuve \\
Heure & 08:00 & Heure de début \\
Durée (minutes) & 120 & Pas de 15 minutes, minimum 15 \\
Coefficient & 1 & Poids dans la moyenne de l'examen (1 à 20) \\
Salle & vide & Par exemple : Salle A1 \\
\bottomrule
\end{tabularx}
\end{center}

\begin{boiteastuce}[Un examen pour plusieurs classes demande plusieurs épreuves]
Si l'examen blanc porte sur trois classes (6ème A, 6ème B, 5ème A en
mathématiques), il faut créer \textbf{trois épreuves}. C'est voulu : chaque classe
a sa propre liste d'élèves et donc ses propres notes.
\end{boiteastuce}

\subsection{Suivre l'avancement}

Un menu déroulant sur la carte fait passer l'examen de \menu{Préparé} à
\menu{En cours}, puis \menu{Terminé} : c'est un repère visuel pour l'équipe.
Chaque ligne d'épreuve porte deux boutons : \bouton{plume} ouvre la saisie des
notes de cette épreuve, \bouton{croix} retire l'épreuve. En bas de la carte, le
total des coefficients est affiché.

\newpage

% =============================================================================
\section{Saisir les notes}

\textbf{Chemin :} \menu{Saisie des notes}.

\subsection{Deux modes, un menu}

En haut à gauche, \champ{Mode de saisie} propose deux choix :

\begin{center}
\renewcommand{\arraystretch}{1.25}
\begin{tabularx}{\textwidth}{@{} >{\raggedright}p{4.6cm} >{\raggedright}p{4.4cm} X@{}}
\toprule
\textbf{Mode} & \textbf{Ce qu'il demande} & \textbf{Pour quoi faire} \\
\midrule
Devoirs et composition & Classe, matière, période & Les notes du trimestre ou du semestre \\
Examen blanc & une épreuve dans une liste & Les notes d'une épreuve blanche \\
\bottomrule
\end{tabularx}
\end{center}

\subsection{La grille en mode période}

\begin{center}
\renewcommand{\arraystretch}{1.3}
\begin{tabular}{@{}lcccc@{}}
\toprule
\textbf{Élève} & \textbf{Devoir 1} & \textbf{Devoir 2} & \textbf{Composition} & \textbf{Moyenne} \\
\midrule
ADJOVI Ama & 12 & 14 & 9 & 10,33 \\
KOUASSI Elom & 15 & 16 & 14 & 14,67 \\
\bottomrule
\end{tabular}
\end{center}

\begin{itemize}
  \item Sous chaque colonne figure une petite case \menu{abs.} : elle marque une
        absence \textbf{sans compter zéro}.
  \item La colonne \textbf{Moyenne} se calcule en direct, ainsi que
        l'appréciation proposée.
  \item Une pastille en haut à droite indique la progression de la saisie.
\end{itemize}

\begin{boiteattention}[Absent n'est pas zéro]
Si vous saisissez \texttt{0} pour un élève absent, sa moyenne s'effondre et son
rang devient injuste. Cochez \menu{abs.} : la mention apparaît sur le bulletin,
et la moyenne est calculée sur les notes réellement obtenues.
\end{boiteattention}

\subsection{Les raccourcis clavier}

\begin{center}
\renewcommand{\arraystretch}{1.25}
\begin{tabularx}{\textwidth}{@{} >{\raggedright}p{3.4cm} X@{}}
\toprule
\textbf{Touche} & \textbf{Effet} \\
\midrule
Entrée & Passe à l'élève suivant, champ présélectionné : on tape la note suivante
sans toucher la souris \\
Ctrl + S & Enregistre toute la grille sans cliquer \\
Maj + Tab & Revient au champ précédent \\
Échap & Ferme la fenêtre ou le menu ouvert \\
\bottomrule
\end{tabularx}
\end{center}

\begin{boiteastuce}[Saisir quarante notes en deux minutes]
Cliquez dans la première case, tapez la note, appuyez sur \textbf{Entrée}, tapez la
suivante, et ainsi de suite. À la fin seulement, \textbf{Ctrl + S}. La main ne
quitte jamais le clavier.
\end{boiteastuce}

\subsection{Période clôturée, et journal de saisie}

\begin{itemize}
  \item Une \textbf{période clôturée} grise la grille et la verrouille pour les
        enseignants. Le chef d'établissement, lui, peut toujours corriger.
  \item Le \textbf{Journal de saisie}, en bas de page et réservé au chef, liste
        qui a enregistré ou modifié chaque note, et quand.
\end{itemize}

\newpage

% =============================================================================
\section{Lire les résultats}

\subsection{Les bulletins d'une classe}

\textbf{Chemin :} \menu{Bulletins}. Trois filtres commandent tout :

\begin{center}
\renewcommand{\arraystretch}{1.25}
\begin{tabularx}{\textwidth}{@{} >{\raggedright}p{3.2cm} X@{}}
\toprule
\textbf{Filtre} & \textbf{Valeurs} \\
\midrule
Classe & la division voulue \\
Source & \emph{Bulletins de période} (T1, T2, T3…) ou \emph{Résultats d'examens blancs} \\
Type de bulletin & \emph{Toutes les notes} (bulletin officiel) ou \emph{Composition seule} \\
\bottomrule
\end{tabularx}
\end{center}

Quatre indicateurs résument la classe : moyenne de classe, taux de réussite
(notes supérieures ou égales à 10), meilleure moyenne, plus faible moyenne. Puis
le tableau, trié par rang :

\begin{center}
\renewcommand{\arraystretch}{1.25}
\begin{tabular}{@{}clcccccl@{}}
\toprule
\textbf{Rang} & \textbf{Élève} & \textbf{Moy.} & \textbf{Compo} & \textbf{Mention} & \textbf{Mat.} & \textbf{Publication} \\
\midrule
1 & KOUASSI Elom & 15,40 & 14 & Bien & 8/8 & Publié \\
2 & ADJOVI Ama & 10,33 & 9 & Passable & 8/8 & Brouillon \\
\bottomrule
\end{tabular}
\end{center}

\begin{itemize}
  \item La colonne \textbf{Mat.} (par exemple \texttt{8/8}) indique combien de
        matières sont évaluées : c'est votre indicateur de saisie terminée.
  \item Si des élèves n'ont pas toutes leurs notes, un bandeau orange prévient :
        \menu{X élève(s) n'ont pas encore de moyenne complète}.
\end{itemize}

\subsection{Les cinq actions}

\begin{center}
\renewcommand{\arraystretch}{1.25}
\begin{tabularx}{\textwidth}{@{} >{\raggedright}p{3.5cm} >{\raggedright}p{4.2cm} X@{}}
\toprule
\textbf{Action} & \textbf{Comment} & \textbf{Effet} \\
\midrule
Appréciation & le crayon, sur la ligne de l'élève & Zone de texte préremplie d'une phrase
cohérente avec la moyenne et le rang ; réécriture libre \\
Publier & \bouton{Publier les bulletins} & Toute la classe devient visible par les élèves ;
un second clic dépublie \\
Imprimer la classe & \bouton{Imprimer la classe} & Tous les bulletins, un par page \\
Un seul élève & l'imprimante, sur sa ligne & Utile pour un élève absent le jour de la
distribution \\
CSV & \bouton{CSV} & Moyennes et rangs, pour Excel ou le conseil de classe \\
\bottomrule
\end{tabularx}
\end{center}

En type \emph{Composition seule}, la moyenne de chaque matière est la note de
composition : le bulletin ne s'appuie que sur celle-ci, et le bouton devient
\bouton{Bulletins de composition}.

\subsection{Les résultats d'un examen blanc}

\textbf{Chemin :} \menu{Examens blancs} $\rightarrow$ icône graphique de l'examen
$\rightarrow$ choisir la classe. Le site ouvre la même page \menu{Bulletins}, en
mode simulation :

\begin{itemize}
  \item il n'y a \textbf{pas} de bouton Publier : rien n'est envoyé aux élèves,
        c'est un document de travail pour la direction ;
  \item le bouton d'impression devient \bouton{Relevé de la classe}.
\end{itemize}

\subsection{Puis les statistiques}

\textbf{Chemin :} \menu{Statistiques}. Moyennes par matière, répartition par
mention, comparaison filles / garçons, taux de réussite : tout se recalcule à
partir des notes, rien n'est saisi à la main.

\newpage

% =============================================================================
\section{La séquence de fin de trimestre}

À afficher au mur du secrétariat. L'ordre compte.

\begin{center}
\renewcommand{\arraystretch}{1.3}
\begin{tabularx}{\textwidth}{@{}c >{\raggedright}p{6.2cm} X@{}}
\toprule
\textbf{Ordre} & \textbf{Action} & \textbf{Où} \\
\midrule
1 & Vérifier que tous les professeurs ont saisi (colonne Mat. complète) & Bulletins \\
2 & Contrôler le journal de saisie & Saisie des notes \\
3 & Écrire les appréciations & Bulletins, crayon \\
4 & Imprimer \textbf{un} bulletin de contrôle et le relire & Bulletins, imprimante d'une ligne \\
5 & Imprimer la classe, ou l'enregistrer en PDF & Bulletins \\
6 & Publier les bulletins & Bulletins, bouton Publier \\
7 & Clôturer la période & Paramètres, Périodes et verrouillage, puis Enregistrer \\
8 & Sauvegarder la base & serveur : \texttt{npm run backup} \\
\bottomrule
\end{tabularx}
\end{center}

\begin{boiteattention}[Ne clôturez jamais avant d'avoir imprimé]
Une période clôturée bloque les enseignants — pas vous — mais autant ne pas
travailler sous contrainte pendant la saisie. Imprimez, distribuez, \emph{puis}
clôturez.
\end{boiteattention}

\section{Annexe — Fiche à cocher}

\subsection*{Création d'une classe}
\begin{itemize}[label=$\square$]
  \item Nom exact de la classe (orthographe retenue une fois pour toutes)
  \item Cycle vérifié : Collège ou Lycée
  \item Professeur titulaire désigné (compte existant ou nom affiché)
  \item Toutes les matières enseignées cochées
  \item Coefficient saisi pour chaque matière, conforme au règlement
\end{itemize}

\subsection*{Rentrée des élèves}
\begin{itemize}[label=$\square$]
  \item Classes créées avant l'import
  \item Liste au format : NOM ; PRÉNOM ; SEXE ; NAISSANCE ; CLASSE ; TUTEUR ; TÉLÉPHONE
  \item Noms de classes identiques à ceux du site
  \item Bilan d'import lu (créés et refusés)
  \item Export CSV téléchargé et conservé
\end{itemize}

\subsection*{Comptes des enseignants}
\begin{itemize}[label=$\square$]
  \item Email de connexion noté (non modifiable ensuite)
  \item Mot de passe dicté et noté
  \item Au moins une case cochée dans la grille d'affectations
  \item Signature importée si le professeur signe les bulletins
\end{itemize}

\subsection*{Examen blanc}
\begin{itemize}[label=$\square$]
  \item Examen créé avec un nom explicite et des dates
  \item Une épreuve par classe et par matière
  \item Durée et coefficient conformes au modèle de l'examen officiel
  \item Notes saisies par épreuve, absences cochées
  \item Relevé de la classe consulté et commenté en conseil
\end{itemize}

\subsection*{Fin de trimestre}
\begin{itemize}[label=$\square$]
  \item Toutes les notes saisies (colonne Mat. complète pour tous)
  \item Journal de saisie consulté
  \item Appréciations écrites
  \item Un bulletin relu à l'écran
  \item Classe imprimée ou enregistrée en PDF
  \item Bulletins publiés
  \item Période clôturée
  \item Base sauvegardée
\end{itemize}

\vfill
\begin{center}
\small\color{gris}
Terminish École — Manuel du chef d'établissement · version 1.1 · octobre 2026\\
CEG BOROME — HIHEATRO (Togo) · https://kodjone.alwaysdata.net
\end{center}

\end{document}
